\section{组织与投资的启示}

\subsection{人才结构与文化}

Burak 特别强调人才密度与迭代文化的重要性：把跨功能 squad 组织成 product + safety + research 的三位一体。每个 squad 内部都需要 prompt engineer、search engineer、ops engineer 与 policy analyst 四类角色，确保 agent pipeline 不仅能跑，还能管。

Anthropic 在 2023 年的经验是：人才密度 > 数量。他们把组织值定义成「谁能马上跑完实验并解释 failure」。这种文化确保 code review 同时覆盖 quality 与 compliance。

\subsection{资本分配与项目组合}

GenAI 投资不是全部押在一个 agent，而是构建多阶段 portfolio：

\begin{itemize}
    \item Stage 1：GPT-style capability research 继续探索 scaling, architecture
    \item Stage 2：Platform / Agent ops 提供 tooling + governance
    \item Stage 3：Productization 把 agent 变成 specific vertical solutions（如 finance assistant）
\end{itemize}

\begin{importantbox}{分阶段 portfolio 的好处}
这样的组合允许团队在 early-stage 上做创新探索，mid-stage 保证 platform reliability，late-stage 担当 commercialization。每个阶段的 metrics / budget 都需要清晰分离，避免能力部门和 product 部门在同一个 sprint 中彼此干扰。
\end{importantbox}

\subsection{本章小结}

组织层面需要建立「人才密度」「跨功能 squad」和「分阶段投资」三大机制，才能在技术快速演进之际保持团队可控。

